\begin{afterword}
\summaryhead

The post builds on earlier posts. Knowledge is mutual information between a mind and its surroundings, and such information cannot be created without breaking the second law of thermodynamics. So a mind that arrives at true beliefs must have processed evidence in a Bayesian way at least a little, and any part of a mind that helps it find truth must have ``at least a little Bayesian structure.''

The author then describes a larger quest. Many cognitive processes do not look Bayesian, and philosophers have argued about them for centuries. On examination, the author reports, each turns out to have Bayesian structure underneath, and the full structure found is always the same, like a Holy Grail. The post's own example is its account of words. Knowing only that a process is Bayesian is a password; the aim is to see how, which the author calls ``Bayes-Sight.''

\responsehead

The post makes two claims of very different strength, and gives an argument only for the weaker one.

The weak claim is that any useful part of a mind has ``at least a little'' Bayesian structure. The thermodynamic argument supports it, and the post states it with its limits: ``at least vaguely Bayesian,'' ``sort-of,'' ``somewhere,'' ``however noisily.'' So qualified, it says little more than its premise. ``Bayesian structure'' is not defined. A part that helps find truth helps make the mind's states carry information about the world, and that is what the argument shows.

The strong claim is that every search ends with ``the same Holy Grail,'' the ``entire'' Grail, and that each cognitive process is Bayesian ``as it always must be.'' For this the post offers the author's experience of repeated discoveries and two linked cases, the earlier posts on words and regularized regression. The second is well chosen: ridge regression does have an exact Bayesian reading. But the procedure the post describes has no outcome that counts against the thesis. When a part looks non-Bayesian, either it turns out Bayesian, or it is judged not to be doing anything useful. A useful part with no Bayesian structure, the one case that would count against the thesis, has no place in the account.

The post is candid about what its argument gives. Knowing that a process is Bayesian without knowing how is, it says, ``a hint at the form an answer would take; but certainly not an answer.'' The thermodynamic argument gives exactly that: some Bayesian structure must exist, but not what it is. The post presents the argument as the start of the quest, not its result. The gap between the two claims is closed only by the reported experience.

Two smaller points. The claim that all the philosophical debate about meaning was ``all along'' Bayesian inference names no philosopher and no position. The view that statistical mechanics is a form of statistics is E.~T. Jaynes's, and the post does not name him.

In short: a sound argument for a weak claim, used to introduce a strong one that rests on the author's reported experience and on a search whose outcomes all count in its favour.
\end{afterword}
